%% ==========================================================================
%%  preamble.tex --- packages, theorem environments, notation macros.
%%
%%  SHARED BY ALL THREE MILESTONE DOCUMENTS:
%%      survey.tex    Milestone 1, literature survey      (31 Aug 2026)
%%      midterm.tex   Milestone 2, per-member report      (12 Sep 2026)
%%      final.tex     Milestone 3, final report           ( 6 Nov 2026)
%%
%%  One preamble and one references.bib for the whole semester, so notation
%%  and citations cannot drift between documents.  If you need a symbol,
%%  add it HERE -- never redefine it inside a section file.
%% ==========================================================================

% ---------- page geometry & fonts -----------------------------------------
\usepackage[letterpaper,margin=1.15in]{geometry}
% Palatino text + matching math.  newpx is built on TeX Gyre Pagella, so it
% needs no URW base-35 fonts -- which a minimal TinyTeX install does not ship.
% If newpx is unavailable, `mathpazo' is the drop-in substitute wherever the
% URW fonts *are* present.
\usepackage[T1]{fontenc}
\usepackage[utf8]{inputenc}
\usepackage{newpxtext}
\usepackage{microtype}

% ---------- mathematics ---------------------------------------------------
% Order matters: amsmath must be loaded before newpxmath.
\usepackage{amsmath,amssymb,amsthm,mathtools}
\usepackage{newpxmath}
\usepackage{stmaryrd}          % \llbracket \rrbracket -- secret-sharing brackets

% ---------- floats, tables, algorithms ------------------------------------
\usepackage{booktabs}
\usepackage{multirow}
\usepackage{longtable}
\usepackage{array}
\usepackage{graphicx}
\graphicspath{{figures/}}
\usepackage{algorithm}
\usepackage[noend]{algpseudocode}
\floatname{algorithm}{Algorithm}
\usepackage{tikz}
\usetikzlibrary{arrows.meta,positioning,fit,backgrounds,calc}

% ---------- colour & links ------------------------------------------------
\usepackage{xcolor}
\definecolor{CiteGreen}{RGB}{0,120,60}
\definecolor{LinkRed}{RGB}{160,20,20}
\definecolor{UrlBlue}{RGB}{20,60,170}
\definecolor{DraftOrange}{RGB}{200,90,0}
\usepackage[
  colorlinks=true,
  citecolor=CiteGreen,
  linkcolor=LinkRed,
  urlcolor=UrlBlue,
  bookmarksnumbered=true
]{hyperref}
\usepackage{cleveref}

% ==========================================================================
%  Theorem environments.  One counter runs through
%  Definition/Theorem/Lemma/Proposition/Observation; Claim has its own.
% ==========================================================================
\theoremstyle{plain}
\newtheorem{theorem}{Theorem}[section]
\newtheorem{lemma}[theorem]{Lemma}
\newtheorem{proposition}[theorem]{Proposition}
\newtheorem{corollary}[theorem]{Corollary}
\newtheorem{definition}[theorem]{Definition}
\newtheorem{observation}[theorem]{Observation}
\newtheorem{conjecture}[theorem]{Conjecture}
\newtheorem{claim}{Claim}[section]

\theoremstyle{definition}
\newtheorem{example}[theorem]{Example}
\newtheorem{remark}[theorem]{Remark}
\theoremstyle{plain}

% NOTE.  cleveref cannot tell these apart --- they share the `theorem'
% counter, so \Cref would call every one of them "Theorem".  Write
% `Remark~\ref{...}', `Example~\ref{...}' explicitly.  \Cref is safe for
% sections, figures and tables.

% Bold, unpunctuated "Proof" head (amsthm defaults to italic "Proof.").
\makeatletter
\renewenvironment{proof}[1][\proofname]{\par
  \pushQED{\qed}%
  \normalfont \topsep6\p@\@plus6\p@\relax
  \trivlist
  \item[\hskip\labelsep\bfseries #1\@addpunct{}]\ignorespaces
}{%
  \popQED\endtrivlist\@endpefalse
}
\makeatother

% ==========================================================================
%  Draft notes.  Flip to \draftfalse for a clean submission copy; nothing
%  else in the document changes.
%
%      \todo{...}    something still to be written
%      \verify{...}  a claim that must be checked against the paper
%      \gap{...}     a place where the survey must argue the gap explicitly
%
%  Before submitting, find what is outstanding:
%      grep -rn '\\todo{\|\\verify{\|\\gap{' sections/ *.tex
% ==========================================================================
\newif\ifdraft
\drafttrue

\newcommand{\draftnote}[1]{%
  \ifdraft{\color{DraftOrange}\sffamily\small[\,#1\,]}\fi}
\newcommand{\todo}[1]{\draftnote{\textbf{TODO:} #1}}
\newcommand{\verify}[1]{\draftnote{\textbf{VERIFY:} #1}}
\newcommand{\gap}[1]{\draftnote{\textbf{GAP-ARGUMENT:} #1}}

% ==========================================================================
%  PROJECT NOTATION.
%
%  Fixed here, once, for all three documents.  Follows [NUDGE] where the two
%  base papers disagree, since its notation is the more standard.  Do not
%  introduce a competing symbol in a section file --- add it here instead,
%  and say so at the weekly sync so nobody duplicates it.
% ==========================================================================

% ---- parties and sharing -------------------------------------------------
\newcommand{\party}[1]{P_{#1}}                 % server P_0, P_1, P_2
\newcommand{\shr}[1]{\llbracket #1 \rrbracket} % replicated / additive sharing
\newcommand{\ring}{\mathcal{R}}                % the ring Z_{2^b}
\newcommand{\Ztwo}[1]{\mathbb{Z}_{2^{#1}}}
\newcommand{\bits}{b}                          % ring bit-length
\newcommand{\fracbits}{t}                      % fixed-point fractional bits
\newcommand{\seclen}{\lambda}                  % PRF seed / security parameter

% ---- the recommendation problem -----------------------------------------
\newcommand{\users}{m}                         % number of users
\newcommand{\items}{n}                         % number of items
\newcommand{\dime}{d}                          % embedding dimension
\newcommand{\iters}{\ell}                      % power-iteration rounds
\newcommand{\topk}{k}                          % recommendations returned
\newcommand{\ratings}{\mathbf{U}}              % user-item rating matrix, m x n
\newcommand{\userembed}{\mathbf{A}}            % user embeddings, m x d (SECRET)
\newcommand{\itemembed}{\mathbf{B}}            % item embeddings, d x n (PUBLIC)
\newcommand{\uvec}[1]{\mathbf{u}^{(#1)}}       % user i's rating vector
\newcommand{\avec}[1]{\mathbf{a}^{(#1)}}       % user i's embedding

% ---- protocols -----------------------------------------------------------
\DeclareMathOperator{\Trunc}{Trunc}
\DeclareMathOperator{\ApproxNormalize}{ApproxNormalize}
\DeclareMathOperator{\ApproxFactor}{ApproxFactor}
\DeclareMathOperator{\SetOrthogonal}{SetOrthogonal}
\DeclareMathOperator{\Normalize}{Normalize}
\DeclareMathOperator{\Gen}{Gen}
\DeclareMathOperator{\Eval}{Eval}
\DeclareMathOperator{\EvalFull}{EvalFull}
\newcommand{\dpf}{\mathrm{DPF}}
\newcommand{\fss}{\mathrm{FSS}}

% ---- shorthand for the two base papers ----------------------------------
%  Use \pirsona and \nudge in prose so the naming stays uniform; the macros
%  cite as well as name, so `\nudge{} publishes $\itemembed$ in the clear'
%  produces both the name and the citation.
\newcommand{\pirsona}{\textsc{Pirsona}\,\cite{pirsona2021}}
\newcommand{\nudge}{\textsc{Nudge}\,\cite{nudge2026}}
\newcommand{\pirsonaplain}{\textsc{Pirsona}}
\newcommand{\nudgeplain}{\textsc{Nudge}}
\newcommand{\sys}{\textsc{OblivRec}}           % our system, working name

% ---- generic -------------------------------------------------------------
\newcommand{\N}{\mathbb{N}}
\newcommand{\R}{\mathbb{R}}
\newcommand{\Z}{\mathbb{Z}}
\newcommand{\set}[1]{\left\{#1\right\}}
\newcommand{\abs}[1]{\left\lvert#1\right\rvert}
\newcommand{\norm}[1]{\left\lVert#1\right\rVert}
\newcommand{\inner}[2]{\left\langle #1, #2 \right\rangle}
\DeclareMathOperator*{\argmax}{arg\,max}
\DeclareMathOperator*{\argmin}{arg\,min}
\DeclareMathOperator{\poly}{poly}
